\section{Limitations and responsible use}
\label{sec:limitations}

The limitations below are not a disclaimer appended for legal comfort. Several
of them --- the daytime restriction, the air-temperature-only scope, the
geographic skew of the evidence --- materially constrain what the tool's
outputs mean, and a user who does not know about them will misread the results.
They are stated at length for that reason.

\subsection{Screening level}

Outputs are indicative estimates produced from simplified typology factors and
from user-supplied or default assumptions. They are not microclimate simulation
outputs, and they must not be presented as predicted temperatures at a
location. The tool has no representation of site geometry, no surface energy
balance, no wind field, and no time dimension beyond the daytime/nighttime
distinction it declines to model.

The practical consequence is that small differences in an Opportunity Score
between two options carry little meaning. The tool distinguishes
a strong opportunity from a weak one; it does not rank two strong ones against
each other with any reliability. Users comparing closely scored options should
treat them as tied and decide on grounds the tool does not represent.

\subsection{Daytime only}
\label{sec:limit-daytime}

All cooling values in this methodology are daytime values. The tool produces no
nighttime estimate.

This is a substantive restriction, not a scoping convenience, because the
evidence indicates that nighttime behaviour differs both in magnitude and in
mechanism. \textcite{keravec2026} reports substantially smaller nighttime
effects --- parks at \qty{1.2}{\degreeCelsius} against
\qty{1.3}{\degreeCelsius} by day, green walls at \qty{1.4}{\degreeCelsius}
against \qty{3.0}{\degreeCelsius}. \textcite{ziter2019} finds canopy cooling
minimal at night while the warming effect of impervious surfaces
\emph{persists} overnight, which means that at night the two mechanisms
represented in the tool's Heat Exposure Score no longer offset one another.

The importance of this is that nighttime heat is strongly associated with
heat-health outcomes: sustained overnight temperatures prevent physiological
recovery, and it is that failure to recover which drives excess mortality
during heat events. A favourable daytime result from this tool should
\textbf{not} be read as evidence of overnight heat relief, and an intervention
selected on the strength of a daytime score may deliver much less of the
health benefit a city is actually seeking.

\begin{keynote}
\noindent\textbf{The evidence retrieved for version \methver{} sharpens this
limitation into a warning for two classes of intervention. They do not merely
cool less at night; they warm.}

\textbf{Water.} \textcite{yao2023b} measures \qty{1.8}{\degreeCelsius} of
nocturnal warming at an urban pond, \textcite{yao2023a} measures a
\SIrange{1.2}{1.3}{\degreeCelsius} nocturnal heat island at urban lakes, and
\textcite{ampatzidis2020} states that blue space may exacerbate the nighttime
urban heat island outright, with sensible cooling partly offset by increased
water-vapour content.

\textbf{Irrigated productive landscapes.} \textcite{cheung2022} predicts an
increase in the daily \emph{minimum} air temperature, as heat stored in wet
soil through the day is released overnight.

In both cases the effect is \textbf{outside} the tool's daytime estimate rather
than netted off against it, and both archetypes carry an output caveat saying
so. A daytime-only tool that quietly omitted this would be reporting only the
favourable half of what its own sources found.
\end{keynote}

A user whose primary concern is nighttime heat should treat this tool's output
as addressing a related but different question.

\subsection{Air temperature only --- thermal comfort is not modelled}
\label{sec:limit-airtemp}

The tool estimates air temperature reduction. It does not quantify thermal
comfort, which is what a person standing in the street actually experiences.

The two diverge most where shade is the dominant mechanism. Shade improves
comfort substantially by reducing radiant load, even where the effect on air
temperature is small, and comfort indices capture this while air temperature
does not. \textcite{zolch2016} reports a \SI{13}{\percent} PET reduction from
trees at pedestrian level, which illustrates the magnitude of what is omitted;
that study also found that placement matters more than quantity, a design
insight this methodology cannot express through an air-temperature envelope.

The direction of the resulting error is known and is worth stating plainly:
\textbf{the tool systematically understates the benefit of shade-dominant
interventions.} Every entry inheriting the street tree canopy, the continuous
shaded corridor, the permeable shaded hardscape, or the vegetated shade
structure delivers more human benefit than its air-temperature figure conveys.
The vegetated shade structure is the starkest case, and it is measured:
\textcite{ouyang2024} reports an air temperature reduction of
\qty{1.42}{\degreeCelsius} alongside a mean radiant temperature reduction of
\qty{15.93}{\degreeCelsius} for the same natural shade --- an order of
magnitude apart. The tool reports the first number and says so.

This bias was not corrected, and the decision not to correct it deserves
explanation. Applying an uplift to shade-dominant typologies would require a
conversion factor between air temperature and comfort that no retrieved source
supplies, and inventing one would trade a known, stated, one-directional bias
for an unknown one concealed inside a number. A stated bias that users can
reason about is preferable to a hidden correction they cannot.

\subsection{The evidence base is geographically uneven}

The underlying literature is dominated by temperate and subtropical Northern
Hemisphere cities. Tropical and arid-climate estimates rest on fewer studies
with wider intervals --- visible in \Cref{fig:climate-dependence}, where the
tropical street-tree estimate carries a \qty{\pm 1.6}{\degreeCelsius} interval
and the arid estimate \qty{\pm 1.4}{\degreeCelsius} on a central value of
\qty{1.7}{\degreeCelsius}, an interval nearly as large as the estimate itself.

The energy sensitivity of \Cref{sec:energy} derives largely from North American
evidence, and its transferability to other building stocks and climates is
untested.

Applications in under-represented contexts should treat outputs as more
uncertain than the confidence rating alone suggests. The confidence model of
\Cref{sec:confidence} measures input completeness; it does not measure how well
the evidence base covers the user's part of the world. A complete assessment of
a site in a tropical city may report high cooling confidence while resting on
thinner evidence than the same assessment in a temperate one. This is a known
gap in the confidence model and a candidate for a future methodology version.

\subsection{An entry's cooling value is its archetype's, not its own}
\label{sec:limit-inheritance}

This is the limitation created by the archetype model of
\Cref{sec:archetypes}, and it is the price of the honesty that model buys.

\num{121} catalogue entries are served by \num{18} cited evidence classes.
Fifteen entries therefore report the street tree canopy's
\SIrange{0.5}{3.0}{\degreeCelsius} and eleven report the green fa\c{c}ade's
\SIrange{0.3}{2.0}{\degreeCelsius}, whatever distinguishes one entry from
another within the class. A suspended-pavement tree pit and a continuous tree
trench are not the same intervention, and the tool does not pretend to
distinguish their cooling.

The alternative was worse. Assigning \num{121} separate envelopes would have
required inventing \num{92} of them, since solution-specific measurement does
not exist at that resolution, and a fabricated envelope is indistinguishable in
a report from a measured one. What the tool does instead is inherit a
\emph{cited} envelope and \textbf{name the evidence class in every result}, so
that a reader can see the resolution of the claim rather than having to infer
it. Two entries reporting the same range because they inherit the same class is
a visible fact about the output, not a hidden assumption inside it.

The consequence for use is direct: \textbf{the tool does not discriminate
between entries within one archetype on cooling}. Two such entries are
separated, if at all, by their suitability conditions, their co-benefit
defaults, and their availability --- not by their temperature range. A user
choosing between them on the strength of a cooling difference is reading a
difference that is not there.

\subsection{Newer and larger published figures are not adopted}
\label{sec:limit-non-adoption}

\Cref{sec:non-adoption} records that \textcite{kumar2024} reports figures
substantially above this library's envelopes for several intervention types ---
green walls at \qty{4.1(42)}{\degreeCelsius}, vegetated balconies at
\qty{3.8(27)}{\degreeCelsius}, botanical gardens at
\qty{5.0(35)}{\degreeCelsius} --- and that they are deliberately not adopted,
because selective adoption would break internal consistency and wholesale
adoption is a full recalibration requiring its own sensitivity analysis.

A reader should understand what that means for the outputs: \textbf{if those
figures are right, this tool understates the corresponding interventions.} The
non-adoption is a defensible methodological choice and it is also a live
disagreement with a recent review, held open rather than resolved. It is
recorded as a candidate for a future methodology version, and a reviewer who
believes the recalibration should happen now is making an argument the project
would need to answer with evidence rather than with precedent.

\subsection{Modelling bias is mitigated, not eliminated}

Much of the synthesised evidence derives from simulation rather than field
measurement, and simulation tends to report larger effects.
\Cref{sec:modelling-bias} describes the downward correction applied through
conservative bound-setting. That correction is a judgment about direction
informed by a single observable ratio; it is not a calibrated adjustment, and
it does not eliminate the bias. Adopted values may still be optimistic.

\subsection{No default costs}

Economic outputs require user-supplied costs (\Cref{sec:costs}). Absent them,
the tool reports no capital cost, no payback, and no cost feasibility, and
excludes the latter from aggregation. A user seeking an investment case from
the tool alone will not get one. This is deliberate, and the reasoning is given
in \Cref{sec:costs}; but it is a real functional limitation and should be
planned for.

\subsection{The weights are contestable}

The aggregation weights, including the equity-forward emphasis quantified in
\Cref{sec:effective-weights}, are value choices. Users who disagree should
adjust them in \texttt{config/weights.yaml} and re-run, rather than
reinterpreting outputs produced under weights they reject. Reinterpretation is
untraceable; reconfiguration is versioned, and the resulting scores carry the
version stamp of the configuration that produced them.

\subsection{Specification gaps at version \methver{}}

The three specification gaps disclosed at version 2026.07.30 --- the
input-to-sub-indicator rules for the NbS Suitability and Equity scores
(\Cref{sec:composite-subscores}), the payback bracket boundaries and
combination rule for cost feasibility (\Cref{eq:cost-feasibility}), and the
enumeration of which input fields feed which confidence block
(\Cref{sec:confidence}) --- were closed at version 2026.08.01 and remain
closed. They are named here because a reader comparing versions should be able
to see what happened to them, not because they are still open.

One gap remains, and it is a data gap rather than a specification gap: the
country emission factor table ships \textbf{empty} (\Cref{sec:ghg}), so
greenhouse-gas outputs are reported as \emph{not calculated} unless a
deployment supplies its own factor. This is the verification policy operating
as designed --- no unverified value ships --- and it is nonetheless a
functional limitation. Closing it requires per-country verification, a version
bump, and a corresponding update to this document in the same change set.

\subsection{Misuse cases}
\label{sec:misuse}

Three specific misuses are anticipated. They are named because each is likely,
each is damaging, and each is easy to commit without intending to.

\subsubsection{Treating outputs as predicted temperatures}

The most likely misuse is quoting a temperature reduction range as a
prediction: ``this project will reduce temperatures by up to
\qty{2.85}{\degreeCelsius}''. It will not. The range is an envelope of what
comparable interventions have delivered in the published literature, narrowed
by a coarse site adjustment. It carries no location, no time of day beyond
``daytime'', no meteorological condition, and no confidence interval in the
statistical sense.

The correct formulation names the metric, the range, and the basis: ``a
screening assessment indicates an indicative daytime pedestrian-level air
temperature reduction in the range \qtyrange{0.5}{2.85}{\degreeCelsius}, based
on published evidence for this intervention type adjusted for site
conditions''. This is longer, and it is what the tool actually supports.

The failure mode is particularly acute where a range is collapsed to its upper
bound for a headline. The upper bound of an envelope describes a well-executed
instance under favourable conditions; it is not the expected outcome.

\subsubsection{Quoting scores without their confidence}

An Opportunity Score of 72.6 at medium confidence and an Opportunity Score of
72.6 at low confidence are different results, and the difference is not
cosmetic --- the second may rest largely on the neutral defaults of
\Cref{sec:unknown} rather than on anything the user knew about the site.

Scores extracted into a slide, a table, or a funding proposal without their
confidence rating and their methodology version have been stripped of the
information needed to judge them. The tool attaches confidence wherever a score
appears, and reports circulated onward should preserve that pairing. Where a
score is quoted, the methodology version should be quoted with it, because the
same site may score differently under a later version.

\subsubsection{Using the tool to justify a decision already taken}

The most consequential misuse is the least visible one. Because the tool is
transparent, configurable, and produces a defensible-looking number, it can be
run after a decision has been made in order to furnish that decision with
evidence --- by selecting the option already chosen, by adjusting weights until
the preferred site ranks first, or by presenting a single assessment without
the alternatives that were considered.

Nothing in the software can prevent this, and no methodology can. What the
design does is leave traces: the methodology version is stamped on every
result, the weights used are recorded in configuration, applied defaults are
itemised, and comparison of options for the same site is a first-class feature
rather than an afterthought. An assessment presented without its alternatives,
or produced under non-default weights that are not disclosed, should be
questioned on those grounds.

The tool is built to compare options and set priorities. Used to ratify a
conclusion reached elsewhere, it lends the appearance of rigour to a process
that had none, which is worse than having no tool at all.

\subsection{Responsible use in summary}

\begin{itemize}
  \item Use the tool to compare options and set priorities, not to justify a
  decision already taken.
  \item Combine results with local knowledge and stakeholder input; the tool
  represents neither.
  \item Validate assumptions through site visits before committing resources.
  \item Re-run assessments as better data becomes available, and expect the
  result to change.
  \item Report the confidence rating and the methodology version alongside any
  score that is quoted.
  \item Name the metric whenever a temperature is reported: daytime,
  pedestrian-level, air temperature.
  \item Carry the evidence class with the number. A cooling range is the
  archetype's, not the entry's (\Cref{sec:limit-inheritance}), and the result
  states which archetype it came from for exactly this reason.
\end{itemize}
